\documentclass[11pt]{article}

% ============================================================
% Packages
% ============================================================
\usepackage[utf8]{inputenc}
\usepackage[margin=1in]{geometry}
\usepackage{amsmath,amssymb,amsthm}
\usepackage{graphicx}
\usepackage{booktabs}
\usepackage{xcolor}
\usepackage{hyperref}
\usepackage{tikz}
\usepackage{enumitem}
\usepackage{pgfplots}

\pgfplotsset{compat=1.18}

% ============================================================
% Document Formatting
% ============================================================
\setlength{\parskip}{0.6em}
\setlength{\parindent}{0pt}

% ============================================================
% List Formatting
% ============================================================
\setlist[enumerate]{
    leftmargin=2em,
    itemsep=0.25em
}

% ============================================================
% Solution Environment
% ============================================================
\newenvironment{solution}
    {\par\noindent\textbf{Solution.}\par}
    {\par}

% ============================================================
% Link Styling
% ============================================================
\hypersetup{
    colorlinks=true,
    linkcolor=blue,
    urlcolor=cyan
}

% ============================================================
% Metadata
% ============================================================
\title{Ordinary Differential Equations: Abridged Notes}
\author{Chris Ritter}
\date{September 17, 2026}

% ============================================================
% Document
% ============================================================

\begin{document}
\maketitle
\begin{center}
\small\emph{These abridged notes were reconstructed afterward because I did not take
direct notes during class. They are not a record of the actual lecture.}
\end{center}
\tableofcontents
\newpage

\section{Topic Summary}
For a homogeneous second-order linear differential equation with
constant coefficients,
\[
ay''+by'+cy=0,\qquad a\ne0,
\]
trying \(y=e^{rt}\) turns the differential equation into the
characteristic equation \(ar^2+br+c=0\). Its roots determine a
fundamental pair of solutions. Distinct real roots give two distinct
exponentials; repeated and complex roots require different forms.
Initial conditions determine the constants in the general solution.

\section{The Characteristic Equation}
If \(r_1\) and \(r_2\) are distinct real roots of
\(ar^2+br+c=0\), then
\[
\boxed{y(t)=C_1e^{r_1t}+C_2e^{r_2t}.}
\]
The two exponentials are linearly independent because their
Wronskian is \((r_2-r_1)e^{(r_1+r_2)t}\ne0\). The form of the
solution changes for a double root or a conjugate pair; those cases
are treated in later notes.

\subsection{Example: Two Distinct Real Roots}
\textbf{Adapted from Boyce et al., 11th ed., \S3.1, Problem 1.}
Solve \(y''+2y'-3y=0\).
\begin{solution}
The characteristic equation factors as
\[
r^2+2r-3=(r-1)(r+3)=0.
\]
The roots are 1 and \(-3\), so
\[
\boxed{y(t)=C_1e^t+C_2e^{-3t}.}
\]
Each exponential satisfies the equation, and their Wronskian
\(-4e^{-2t}\) is nonzero.
\end{solution}

\subsection{Example: An Initial-Value Problem and Minimum}
\textbf{Adapted from Boyce et al., 11th ed., \S3.1, Problem 14.}
Solve \(y''-y=0\), \(y(0)=5/4\), \(y'(0)=-3/4\); find its minimum
for \(0\leq t\leq2\).
\begin{solution}
The characteristic roots are \(1\) and \(-1\), so
\(y=C_1e^t+C_2e^{-t}\). At \(t=0\), the data give
\[
C_1+C_2=\frac54,\qquad C_1-C_2=-\frac34.
\]
Thus \(C_1=1/4\), \(C_2=1\), and
\[
y(t)=\frac14e^t+e^{-t}.
\]
Differentiating gives \(y'=e^t/4-e^{-t}\). Its only critical point
is \(t=\ln2\), which lies in \([0,2]\). Since
\(y''=e^t/4+e^{-t}>0\), this point is a minimum, with
\[
\boxed{y(\ln2)=1.}
\]
\end{solution}
\end{document}
